% CHƯƠNG 2 - Kiến thức nền tảng cần để hiểu phương pháp
\chapter{Kiến thức nền tảng}
\label{chuong:nentang}

\section{Mạng nơ-ron hồi quy}
\label{sec:rnn}

Mạng nơ-ron hồi quy (Recurrent Neural Network, RNN) là họ mô hình học sâu xử lý dữ liệu dạng
chuỗi bằng cách duy trì một trạng thái ẩn (hidden state) được cập nhật tuần tự qua từng bước
thời gian, cho phép mô hình ``nhớ'' thông tin từ các bước trước \cite{elman1990rnn}. Đây là
kiến trúc được bài báo gốc dùng cho cả ba mô hình (teacher, baseline, student) trên bài toán
Mackey--Glass.

Với đầu vào $x_t$ tại bước thời gian $t$, RNN cập nhật trạng thái ẩn $h_t$ theo công thức
\begin{equation}
	h_t = \tanh\!\left(W_{xh} x_t + W_{hh} h_{t-1} + b_h\right),
	\label{eq:rnn}
\end{equation}
trong đó $W_{xh}$, $W_{hh}$ là các ma trận trọng số và $b_h$ là độ lệch (bias); $h_t$ vừa là
trạng thái được truyền sang bước thời gian tiếp theo, vừa được dùng để tính đầu ra tại bước
$t$. Nhược điểm được biết đến của RNN là hiện tượng \textit{tiêu biến/bùng nổ gradient}
(vanishing/exploding gradient): khi lan truyền ngược qua nhiều bước thời gian, gradient phải đi
qua tích của nhiều ma trận $W_{hh}$ liên tiếp, khiến gradient dễ tiến về không hoặc tăng vô
hạn, làm mô hình khó học được các quan hệ phụ thuộc xa trong chuỗi.

\section{Chưng cất tri thức}
\label{sec:kd}

Chưng cất tri thức (Knowledge Distillation, KD), đề xuất bởi Hinton và cộng sự
\cite{hinton2015distillation}, là kỹ thuật huấn luyện một mô hình (student) để bắt chước đầu
ra của một mô hình khác (teacher) đã huấn luyện trước, nhằm truyền ``tri thức'' từ teacher
sang student. Ý tưởng cốt lõi là dùng \textit{soft targets} --- phân phối xác suất đầy đủ trên
mọi lớp mà teacher dự đoán --- thay vì chỉ dùng nhãn \textit{cứng} (hard label, một-hot) như
khi huấn luyện thông thường, vì phân phối đầy đủ chứa thêm thông tin về mức độ tương tự giữa
các lớp (``dark knowledge''). Đây chính là kỹ thuật được bài báo gốc dùng làm cầu nối giữa
teacher và student.

Để so sánh hai phân phối xác suất rời rạc $p$ và $q$, KD dùng độ phân kỳ Kullback--Leibler:
\begin{equation}
	D_{\mathrm{KL}}(p \,\|\, q) = \sum_{k} p_k \log \frac{p_k}{q_k},
	\label{eq:kl}
\end{equation}
đo lượng thông tin bị mất khi dùng $q$ để xấp xỉ $p$; $D_{\mathrm{KL}}(p\|q) \geq 0$, bằng 0
khi và chỉ khi $p = q$. Trước khi tính KL, logits của cả teacher và student thường được làm
``mềm'' bằng một tham số nhiệt độ (temperature) $T > 0$ qua hàm softmax có nhiệt độ:
\begin{equation}
	\operatorname{softmax}_T(z)_k = \frac{\exp(z_k / T)}{\sum_j \exp(z_j / T)}.
	\label{eq:softmax-t}
\end{equation}
Khi $T$ lớn, phân phối xác suất trở nên dàn đều hơn, giúp student học được quan hệ tương đối
giữa các lớp có xác suất thấp, chứ không chỉ riêng lớp có xác suất cao nhất. Vai trò cụ thể
của $T$ và của tham số $\alpha$ trong hàm mất mát của bài báo gốc được trình bày chi tiết ở
Chương~\ref{chuong:phuongphap}.

\section{Hệ Mackey--Glass}
\label{sec:mackeyglass}

Hệ Mackey--Glass \cite{mackeyglass1977} là một phương trình vi phân có trễ (delay differential
equation), ban đầu được đề xuất để mô phỏng một số quá trình sinh lý có phản hồi trễ, với dạng
\begin{equation}
	\frac{dy}{dt} = \beta \frac{y(t-\tau)}{1 + y(t-\tau)^n} - \gamma\, y(t),
	\label{eq:mackeyglass}
\end{equation}
trong đó $\tau$ là độ trễ (thường dùng $\tau = 17$), còn $\beta$, $\gamma$, $n$ là các tham số
điều khiển hình dạng của phương trình. Với các giá trị tham số phù hợp, nghiệm của phương
trình \eqref{eq:mackeyglass} có tính \textit{hỗn loạn} (chaotic): quỹ đạo không lặp lại theo
chu kỳ cố định và rất nhạy cảm với điều kiện ban đầu, khiến chuỗi do hệ này sinh ra trở thành
một benchmark chuẩn và khó cho các bài toán dự báo chuỗi thời gian phi tuyến, đặc biệt khi tầm
dự báo (forecast horizon) tăng lên. Đây là bộ dữ liệu đầu tiên trong hai bộ dữ liệu được tái
hiện trong báo cáo này.

\section{Tín hiệu điện não đồ và biến đổi STFT}
\label{sec:eeg}

Điện não đồ (electroencephalography, EEG) ghi lại hoạt động điện của não qua các điện cực đặt
trên da đầu (scalp EEG). Tín hiệu EEG trong cơn động kinh (ictal) thường có đặc trưng tần số
khác biệt rõ rệt so với tín hiệu bình thường (interictal), nên một hướng tiếp cận phổ biến là
biến đổi tín hiệu từ miền thời gian sang miền tần số trước khi đưa vào mô hình học máy --- đây
là cách bài báo gốc xử lý bộ dữ liệu thứ hai được tái hiện trong báo cáo này.

Biến đổi Fourier thời gian ngắn (Short-Time Fourier Transform, STFT) chia tín hiệu thành các
cửa sổ ngắn chồng lắp hoặc không chồng lắp, áp dụng biến đổi Fourier rời rạc lên mỗi cửa sổ, từ
đó tạo ra một \textit{spectrogram} --- biểu diễn năng lượng của tín hiệu theo cả hai chiều
thời gian và tần số. Spectrogram của nhiều kênh EEG có thể được xem như một ``ảnh'' đa chiều
(thời gian × tần số × kênh), cho phép áp dụng các kiến trúc học sâu thiết kế cho ảnh hoặc chuỗi
đa chiều, chẳng hạn kết hợp mạng tích chập (CNN) với mạng hồi quy.

\section{Các độ đo đánh giá}
\label{sec:metrics}

Báo cáo này dùng hai độ đo đánh giá tương ứng với hai loại bài toán.

\subsection{Sai số bình phương trung bình MSE}

Với bài toán hồi quy (dự báo giá trị chuỗi Mackey--Glass), sai số bình phương trung bình
(Mean Squared Error, MSE) được dùng để đánh giá:
\begin{equation}
	\text{MSE} = \frac{1}{N} \sum_{i=1}^{N} \left(y_i - \hat{y}_i\right)^2,
	\label{eq:mse}
\end{equation}
trong đó $y_i$ là giá trị thực tế, $\hat{y}_i$ là giá trị dự báo, và $N$ là số mẫu đánh giá.
MSE càng nhỏ, dự báo càng chính xác.

\subsection{Diện tích dưới đường cong ROC}

Với bài toán phân loại sự kiện (phát hiện/dự báo cơn động kinh), độ đo chính là diện tích dưới
đường cong ROC (Area Under the ROC Curve, AUC-ROC). Đường cong ROC biểu diễn tỉ lệ dương thật
(true positive rate) theo tỉ lệ dương giả (false positive rate) khi thay đổi ngưỡng phân loại;
AUC là diện tích dưới đường cong này, nhận giá trị trong $[0,1]$ \cite{fawcett2006roc}. AUC
bằng 0,5 tương ứng với một bộ phân loại ngẫu nhiên (không có khả năng phân biệt lớp); AUC càng
gần 1, mô hình phân biệt hai lớp (ictal/interictal) càng tốt. Độ đo này được chọn (thay vì độ
chính xác, accuracy) một phần vì dữ liệu ictal và interictal mất cân bằng nghiêm trọng về số
lượng, khiến accuracy dễ gây hiểu nhầm.
